\documentclass[a4paper,11pt,oneside]{article}
\usepackage{a4wide}                     % Iets meer tekst op een bladzijde
\usepackage[english]{babel} 
\usepackage{amsmath}                    % Uitgebreide wiskundige mogelijkheden
\usepackage{amssymb}                    % Voor speciale symbolen zoals de verzameling Z, R...
\usepackage{url}                        % Om url's te verwerken
\usepackage{graphicx}                   % Om figuren te kunnen verwerken
\usepackage[small,bf,hang]{caption}    % Om de captions wat te verbeteren
\usepackage{xspace}                     % Magische spaties na een commando
\usepackage[latin1]{inputenc}           % Om niet ascii karakters rechtstreeks te kunnen typen
\usepackage{float}                      % Om nieuwe float environments aan te maken. Ook optie H!
\usepackage{flafter}                    % Opdat floats niet zouden voorsteken
\usepackage{listings}                   % Voor het weergeven van letterlijke text en codelistings
\usepackage{marvosym}                   % Om het euro symbool te krijgen
\usepackage{textcomp}                   % Voor onder andere graden celsius
\usepackage[left=2.5cm,top=1.4cm,right=1.5cm,bottom=25mm,nohead]{geometry}   % marges instellen
\usepackage{fancyhdr}                   % Voor fancy headers en footers.
\pagestyle{fancy}                       % De bladzijdestijl
\usepackage{graphics}                   % Om figuren te verwerken.
\usepackage{lastpage}										% Om aan de laatste pagina te kunnen.
\usepackage{csquotes}
\usepackage{biblatex}
\addbibresource{sources.bib}
\newcommand{\npar}{\vspace{0.5cm} \par \noindent}      % witruimte aan paragrafen

% begin van het document
\begin{document}


% Header (titel en datum aanpassen)
\fancyhf{}                              % Resetten fancy settings.
\renewcommand{\headrulewidth}{0pt}      % Geen lijn onder de header.
\hspace{-7mm}\begin{tabular}{p{50mm}c}
\includegraphics[width=40mm]{logo}& 
	\begin{tabular*}{120mm}[b]{c}
	\huge{Emission spectroscopy:}\\
	\huge{Gas temperature}\\
	\vspace{1cm}\\
	\Large{Arthur Adriaens \& Atakan Atmaca} \\
	\end{tabular*}
\end{tabular}
\vspace{1cm}

% De inhoud
\section{Typical Discharge Spectrum}
\printbibliography
\end{document}
